\chapter{Introduction}
% SMC background
% SQMC and its benefits
% what makes SQMC parallelizable 
% Football model and why SQMC works in this context
% Our contribution two parts - 1) GPU implementation of SQMC in cuthbert and 2) Modelling football matches with SQMC
% asynchronous, sparse local observations

Sequential Quasi-Monte Carlo (SQMC) is an extension of the Sequential Monte Carlo (SMC) algorithm, replacing the Monte Carlo proposal with deterministic propagation through low-discrepancy point sets. The error rate of SQMC is smaller than the Monte Carlo rate of $\mathcal{O}_P (N^{-1/2})$, and we can expect the algorithm to reach faster convergence for a smaller class of functions \citep{chopingerber2015sqmc}. Furthermore, recent work suggests that, for a fixed particle count, the error of a standard SMC exceeds any fixed threshold at infinitely many time instants with probability equal to one, while SQMC's error vanishes uniformly in time as $N$ grows \citep{gerber2026}. This safety guarantee, on a linear-Gaussian filtering problem, provides a justification for the accuracy advantage of SQMC over SMC in long-horizon problems. 

However, the computational cost of SQMC is $\mathcal{O}(N \log N)$, where $N$ is the number of simulations at each iteration, compared with SMC, which is $\mathcal{O}(N)$. This overhead stems from two additional steps in SQMC---sorting the particles along a Hilbert curve before resampling and applying the inverse-CDF transform to the quasi-random points. This limits the feasibility of SQMC in experiments that require a large number of particles.

Another issue with SQMC is that its performance advantage over SMC decreases in increasing dimensions. The most likely factor is the curse of dimensionality that affects SMC \citep{chopingerber2017sqmcintroduction}. Importance sampling, a key component of the SMC algorithm, suffers from weight degeneracy in high dimensions: the larger the dimension, the greater the discrepancy between the proposal and the target distribution, so that a small fraction of particles carry non-negligible weight at each step \citep{snyderbengtssonbickel2008,bickellibengtsson2008}.

This provides the motivation for this dissertation. In the second chapter, we cover the background of SMC, QMC and SQMC. We then describe our implementation of the SQMC algorithm that is computationally efficient by leveraging parallel accelerators (GPU, TPU and multi-core CPU) using JAX, enabling us to scale SQMC effectively. In the third chapter, we apply SQMC to a specific problem of modelling football matches. Although the latent state tracks the attacking and defensive strengths of all teams, each observation---a single match---depends only on the two teams involved. The likelihood is therefore sparse, and the 'effective dimension' of the problem---the number of coordinates that materially influence the integrand \citep{owen2003}---is far smaller than the nominal state dimension. By modelling the transition distribution as linear-Gaussian, we can exploit this structure by Rao-Blackwellisation: we integrate out the non-playing teams analytically and run SQMC on the smaller block of playing teams. This mitigates the weight degeneracy described above and allows us to evaluate the performance of SQMC in this context. We conclude with a discussion of our results and directions for further research.

This work is part of a contribution to the open-source package `cuthbert`, a library for efficient particle filter algorithms in JAX. All code and experiments are available at \url{https://github.com/ryantjx/sqmc-cuthbert} and \url{https://github.com/state-space-models/cuthbert}.